\scanpage{80}
\mspage{105}
\begin{flushright}n° 371\end{flushright}

Le cas le plus important est celui où $M$ est un $A$-module \uline{libre de dimension finie} $n$ ; on identifie alors canoniquement $\handwritten{\bigwedge}(M^{\handwritten{*}})$ et $(\handwritten{\bigwedge}(M))^{\handwritten{\unsure{*}}\mathrm{gr}}$ au moyen de l'isomorphisme canonique $\handwritten{\lambda}_M$ ; un multivecteur (resp. $p$-vecteur) sur $M^{\handwritten{*}}$ est donc identifié à une forme linéaire sur $\handwritten{\bigwedge}(M)$ (resp. $\overset{p}{\handwritten{\bigwedge}}M$) et appelé une \uline{multiforme} (resp. une $p$-\uline{forme}). \handwritten{*}Avec cette identification, la formule (22) s'écrit
\[
x_1\handwritten{\wedge}x_2\handwritten{\wedge}\ldots\handwritten{\wedge}x_n\ ,\ x_1^{\handwritten{*}}\handwritten{\wedge}x_2^{\handwritten{*}}\handwritten{\wedge}\ldots\handwritten{\wedge}x_n^{\handwritten{*}}=\det(\handwritten{\langle}x_{\unsure{i|1}}\ ,\ x_j^{\handwritten{*}}\handwritten{\rangle}).\ \handwritten{*}
\tag*{(27)}
\]

Si $(e_{\unsure{i|1}})_{1\handwritten{\leq}i\handwritten{\leq}n}$ est une base de $M$, $(e_i^{\handwritten{*}})$ la base duale dans $M^{\handwritten{*}}$, les éléments
\[
e_H^{\handwritten{*}}=e_{i_1}^{\handwritten{*}}\handwritten{\wedge}e_{i_2}^{\handwritten{\unsure{*}}}\handwritten{\wedge}\ldots\handwritten{\wedge}e_{i_p}^{\handwritten{*}}
\tag*{(28)}
\]
(avec les notations de la prop.\unsure{3|8}) forment une base de $\handwritten{\bigwedge}(M^{\handwritten{*}})$ lorsque $H$ parcourt les parties de $\unsure{\handwritten{[\![}1,n\handwritten{]\!]}|\handwritten{[}1,n\handwritten{]}}$ ; cette base n'est autre que la \uline{base duale} de la base $(e_H)$ de $\handwritten{\bigwedge}(M)$, car il résulte de (27) que si $H$, $K$ sont deux parties distinctes de $p$ éléments de $\unsure{\handwritten{[\![}1,n\handwritten{]\!]}|\handwritten{[}1,n\handwritten{]}}$, $(i_h)_{1\handwritten{\leq}h\handwritten{\leq}p}$, $(j_h)_{1\handwritten{\leq}h\handwritten{\leq}p}$ les suites strictement croissantes des éléments de $H$ et de $K$ respectivement, on a
\[
\handwritten{\langle}e_{i_1}\handwritten{\wedge}e_{i_2}\handwritten{\wedge}\ldots\handwritten{\wedge}e_{i_p}\ ,\ e_{j_1}^{\handwritten{*}}\handwritten{\wedge}e_{j_2}^{\handwritten{*}}\handwritten{\wedge}\ldots\handwritten{\wedge}e_{j_p}^{\handwritten{*}}\handwritten{\rangle}=0
\]
puisqu'il existe un indice $k$ \unsure{tel|tul} que $j_k\handwritten{\notin}H$, donc $\handwritten{\langle}e_{i_h}\ ,\ e_{j_k}^{\handwritten{*}}\handwritten{\rangle}=0$ quel que soit $h$.

Si $M$ et $N$ sont deux $A$-modules libres de type fini, $u:M\handwritten{\to}N$ une application $A$-linéaire, ${}^{t}(\handwritten{\bigwedge}(u))$ s'identifie canoniquement à $\handwritten{\bigwedge}({}^{t}u)$ en vertu de la commutativité du diagramme (25) ; il en résulte que si $u$ est \uline{bijectif}, le contragrédient ${}^{t}(\handwritten{\bigwedge}(u))^{-1}$ s'identifie à $\handwritten{\bigwedge}({}^{t}u^{-1})$.

\uline{Remarque}.- A) \uline{Rabiot qui doit remonter au} \S3 (\uline{Algèbres graduées}), \uline{où il a été omis par inadvertance}.

Soient $A$ un anneau commutatif gradué de type $Z$ à degrés positifs, $E$ une $A$-algèbre graduée de type $Z$ à degrés positifs, $(A_{\unsure{1|n}})$ et $(E_n)$ les graduations respectives de $A$ et de $E$. En tant que groupe additif, $E$ est un sous-groupe du \uline{produit}
% archive sheet number: 79 (preceded by a crossed-out unreadable archive number)
